\heading{数据结构与算法A-17秋-期中}
\noteinfo{题目来源为真题扫描件转录。客观题答案和部分参考解答已整理在解答中；原扫描中考场纪律、装订线等非题目信息已尽量删去。本次整理提供 C 语言版与 Python 语言版；除程序代码语言外，两版题目内容保持一致。}

\begin{question}
一、 选择与填空（每空2 分，共24 分）\\
1.\\
下面函数的时间复杂度是（    ）\\
def recursive(int n, int m, int k)\{\\
if (n <= 0)\\
print(m, k)\\
else \{\\
recursive(n-1, m+1, k);\\
recursive(n-1, m, k+1);\\
\}\\
\}\\
A. O(n*m*k)        B. O(n\textasciicircum{}2*m\textasciicircum{}2)\\
B.\\
O(2\textasciicircum{}n)           D. O(n!)\\
2.\\
完成在双循环链表结点p 之后插入s 的操作为(       ):\\
A. p.next.prev=s; s.prev=p; s.next=p.next; p.next=s;\\
B. p.next.prev=s; p.next=s; s.prev=p; s.next=p.next;\\
C. s.next=p.next; s.prev=p; p.next.prev=s; p.next=s;\\
D. s.prev=p; s.next=p.next; s.prev.next=s; s.next.prev=s;\\
3.\\
设栈S 和队列Q 初始状态为空，元素e1，e2，e3，e4，e5，e6 依次通过栈S，一个元素出\\
栈后即进队列Q，若6 个元素出队序列是e2，e4，e3，e6，e5，e1，则栈S 的容量至少是\\
(       )\\
A.2;\\
B.3;\\
C.4;\\
D.6\\
4.\\
设循环队列的容量为40 （序号从0 到39） ，现经过一系列的入队和出队运算后： 1）\\
front=12，rear=19；2）front=19，rear=12；在这两种情况下，循环队列中的元素个数分别为\\
和        。\\
5.\\
一个n 位的字符串共有                    个子串。\\
6.\\
已知二叉树有n 个结点 (n > 0)，则度为2 的结点最多有            个。\\
7.\\
用数组存储一个有50 个元素的最小值堆。已知堆中没有重复元素。给出最大元素的下标\\
可能的取值范围为           到          。\\
8.\\
假设用于通信的电文由5 个字符组成。已知其中一个字符的Huffman 编码为1101，则\\
该Huffman 编码树的平均编码长度为           。\\
9.\\
假设先根次序遍历某棵树的结点次序为GABCDIJEFH，后根次序遍历该树的结点次序为\\
BADIJCFEHG，那么I 结点的父结点是            。\\
10. 一个拥有144 个结点的完全三叉树，现在从倒数第二层上任取一个结点，该结点为叶结点\\
\\
二、 辨析与简答(共24 分)\\
1.\\
（6 分）现有一个由单链表和循环单链表构成的特殊链表，单链表的头元素是head，末尾\\
元素为tail，head.….tail 即是一条普通的链表，同时tail 元素还是另一条循环单链表的头\\
元素。如A.B.C.D.E.F.G.H.E 就是一个特殊链表，其中head 为A，tail 为E。\\
现在你只知道head，但是你想知道这个循环链表中有多少元素，但同时你的剩余空间十分\\
的小，所以你想用尽可能小的空间来解决问题，请问你会怎么做？（依据占用空间给分）\\
2.\\
（6 分）假设以I 和O 分别表示入栈和出栈操作，栈的初始和终态均为空，入栈和出栈的\\
操作序列可表示为仅由I 和O 组成的序列，称可以操作的序列为合法序列，否则为非法序\\
列。如 IOIIOIOO 合法，IOOIOIIO 和IIIOIOIO 为非法。请给出一个算法，判断所给操作\\
序列是否合法。\\
3.\\
（6 分）给定一棵用数组存储的完全二叉树\{10, 5, 12, 3, 2, 1, 8\}\\
（1） 用筛选建堆法将该完全二叉树调整为最大值堆。画出调整后得到的最大值堆，并说明\\
在调整过程中都依次交换了哪些元素。（注：画堆只需要写出层次遍历序列结果即可）\\
（2） 将6、7、14 按顺序插入（1）得到的最大值堆，堆的空间足够大，画出插入14 后得\\
到的最大值堆，并说明在插入14 的过程中都依次交换了哪些元素。\\
（3） 对（2）得到的堆进行3 次deleteMax，画出第3 次deleteMax 后得到的堆，并说明在\\
第3 次deleteMax 的过程中都依次交换了哪些元素。\\
4.\\
（6 分）对下列15 个等价对进行合并，给出所得等价类树的图示。在初始情况下，集合中\\
的每个元素分别在独立的等价类中。使用重量权衡合并规则，合并时子树结点少的并入结\\
点数多的那棵（多的那个作为新树根，少的那个根作为新根的直接子结点）；若两棵树规模\\
同样大，则把根值较大的并入根值较小（新树根取值小的）。同时，采用路径压缩优化。\\
(0,2) (1,2) (3,4) (3,1) (3,5) (9,11) (12,14) (12,9) (4,14) (6,7) (8,10) (8,7) (7,11) (10,15) (10,13)\\
三、 算法填空(每空3 分，共24 分)\\
1.\\
下面的算法利用一个栈将一给定的序列从小到大排序。长度为n的序列由1, 2, …, n这n个连\\
续的正整数组成。请补全下面的代码段，使其可以判断给定的序列是否可以利用一个栈进\\
行排序，如果可以，输出相应的操作。\\
例如：序列4 3 1 2，此序列可以排序，输出：push push push pop push pop pop pop。\\
template <class T>                            // 栈的元素类型为T\\
class Stack \{\\
public:\\
def clear();\\
// 变为空栈\\
def push(const T item);\\
//  item 入栈\\
T pop();\\
// 返回栈顶内容并弹出\\
T top();\\
// 返回栈顶内容但不弹出\\
isEmpty();\\
// 若栈已空返回真\\
isFull();\\
// 若栈已满返回真\\
\};\\
\\
\\
Stack<int> s;\\
//  s 为栈\\
int sequence[maxLen],  len;                 // sequence 为待排序序列， len 为序列长度\\
islegal()  \{\\
// 判断序列是否可以排序\\
int i, j, k;\\
for (i = 0; i < len; i++) \{\\
for (j = i+1; j < len; j++) \{\\
for (k = j+1; k < len; k++) \{\\
if (  //填空 1   )\\
return False;\\
\}\\
\}\\
\}\\
return True;\\
\}\\
def transform()  \{\\
// 输出转换操作\\
int cur = 1,  i = 0;\\
while  (  //填空2  ) \{\\
while (s.isEmpty()  or  //填空3 ) \{\\
//填空4\\
cout << "push ";\\
if (cur == s.top()) break;\\
\}\\
s.pop();\\
print("pop", end=" ")\\
++cur;\\
\}\\
\}\\
2.\\
下面的算法将一个用带右链的先根次序法表示的森林转换为用带度数的后根次序法表示。\\
请利用题目给出的树结点ADT和栈ADT，填充空格，使其成为完整的算法。\\
template<T>      // 数据结构\\
class RlinkTreeNode \{\\
int ltag;                // 左标记\\
T info;\\
RlinkTreeNode<T> *rlink; // 右链\\
\}\\
template<T>\\
class PostTreeNode \{\\
T info;\\
int degree;              // 度数\\
\}\\
// 算法描述：递归遍历用带右链的先根次序法表示的森林，同时计算出度数并转成后根次\\
序森林；设带右链的先根次序法表示的森林为数组\\
RlinkTreeNode RlinkTree[n];\\
// 带度数的后根次序法表示的森林为数组\\
PostTreeNode PostTree[n];\\
int count = 0;                      // 计数器\\
int convert(PostTreeNode * node) \{ // 返回值表示所有右兄弟的数量（包括自己）\\
int tmp;\\
if (1 == node.ltag) \{           // 若没有左子结点，将该结点输出到PostTree\\
// 填空5\\
PostTree[count].degree = 0; // 度数必然是0\\
count++;\\
\}\\
else \{                        // 有左子结点则压栈\\
tmp = convert(node+1);     // 访问第一个子结点，并返回度数\\
PostTree[count].info = node.info;\\
// 填空6\\
count++;\\
\}\\
if (       // 填空7         ) \{  // 有右兄弟(，访问并返回右边兄弟的数量+1)\\
return convert(node.rlink) + 1;\\
\}\\
else \{                       // 没有右兄弟\\
// 填空8\\
\}\\
\}\\
四、 算法设计与实现 (14 分)\\
1.\\
（8 分） 编写算法伪码，对给定的字符串str，返回其最长重复子串及其下标位置。如\\
str=”abcdaccdac”，则子串”cdac”是str 的最长重复子串，下标为2。\\
（重复子串允许重叠）\\
2.\\
（6 分）给出一个算法，求一棵树的树高。可以写出伪代码，需要相应的注释，或者写\\
出详细算法思路。对时间复杂度无具体要求。\\
五、 分析证明题 (共 14 分)\\
1.\\
（8 分）二叉树的内部路径长度是指所有节点的深度的和。假设将N 个互不相同的随机\\
元素插入一棵空的二叉搜索树。请证明得到的二叉搜索树的内部路径长度期望为\\
O(NlogN)。\\
2.\\
（6 分）证明按先根次序周游树获得的序列与其对应二叉树的前序周游获得的序列相同。\\
\\
\\
\\
\

\begin{solution}
本题为整卷转录型资料：选择题、填空题的参考答案以及简答、算法设计题的参考做法已统一放在本解答中，题面不保留原扫描夹带的答案标注。对扫描不清处，保留了原转录中可辨认的信息，后续可继续按原 PDF 图示补充图片。

\medskip
\noindent 原转录中夹带在题面里的参考答案/解答摘录如下，现已移入解答：
本卷包含选择填空、简答、算法设计等题型。下面保留按扫描件整理的题面；其中客观题参考答案已填入空格或括号，主观题给出参考思路或代码。


\end{solution}
\end{question}
